\documentclass[11pt]{article}
\usepackage{fontspec}
\setmainfont{Latin Modern Roman}
\usepackage{amsmath,amssymb,mathtools,unicode-math}
\setmathfont{Latin Modern Math}
\usepackage{xurl,hyperref}
\hypersetup{hidelinks,pdftitle={Problem statement}}
\newcommand{\sref}[2]{\hyperlink{source:#1:#2}{[S#2]}}
\newcommand{\cataloguescope}[1]{\paragraph{Mathematical question.}#1}
\begin{document}
\section{Research Problems in Function Theory — Problem 5.26}\label{statement}

\subsection{Definitions and mathematical statement}

\paragraph{Shared notation.}$\mathbb N=\{1,2,\ldots\}$, and $\mathbb Z,\mathbb Q,\mathbb R,\mathbb C$ denote the integers, rationals, reals and complex numbers. Any other conventions are specified locally.


Write $\mathbb D=\{z\in\mathbb C:|z|<1\}$ and $\mathbb T=\partial\mathbb D$. All Taylor coefficients below are taken at $0$.The Bergman space $A^2(\mathbb D)$ is the Hilbert space of holomorphic functions with
\[\|f\|_{A^2}^2=\frac1\pi\int_{\mathbb D}|f(z)|^2\,dA(z)=\sum_{n=0}^\infty\frac{|a_n|^2}{n+1}<\infty,
\]
where $dA$ is planar area. The condition is equivalent to $\sum_{n\ge1}|a_n|^2/n<\infty$ used in the source. A linear subspace $L\subset A^2(\mathbb D)$ is invariant under the Bergman shift $M_z$ if $zf\in L$ for every $f\in L$. A closed subspace means closed for this Hilbert norm.
\paragraph{Statement recorded in the supplied source.}
Describe the invariant subspaces of $A^2(\mathbb D)$, in particular classify all its closed $M_z$-invariant subspaces. Give a necessary and sufficient structural description, analogous in purpose to the Hardy-space invariant-subspace theorem. The zero subspace and the whole space are included. \sref{statement}{2}
\subsection{Short English statement}
Classify the invariant subspaces of the Bergman shift, including a structural description of every closed invariant subspace.
\subsection{Sources}
\begin{description}
\item[\hypertarget{source:statement:1}{S1}] ULAM.AI and the UnsolvedMath Contributors, UnsolvedMath: A Curated Collection of Open Mathematics Problems, catalogue entry. CC BY 4.0. \url{https://huggingface.co/datasets/ulamai/UnsolvedMath}.
\item[\hypertarget{source:statement:2}{S2}] W. K. Hayman and E. F. Lingham, Research Problems in Function Theory (2018), arXiv:1809.07200, Chapter 5, Problem 5.26, its update and the preceding shared notation. \url{https://arxiv.org/abs/1809.07200}.
\end{description}
\label{end:statement}
\end{document}
